\section{Introduction}\label{sec:introduction}

Penrose's mass--area inequality connects the geometry of initial data to
the expected evolution of black holes under weak cosmic censorship
\cite{Penrose1973}. In the absence of charge, its characteristic form is
$m\ge\sqrt{A/(16\pi)}$. In time-symmetric data, where the scalar
curvature is nonnegative and horizons are minimal surfaces,
Huisken and Ilmanen proved the three-dimensional bound using the area of
any connected component of an outermost minimal boundary, and Bray proved
it using the total area of a possibly disconnected outer minimizing
minimal boundary \cite{HuiskenIlmanen2001,Bray2001}. The passage to an
arbitrary second fundamental form requires an area quantity attached to
the original spacelike exterior. The minimum-area enclosure in the
spacetime formulation of Bray and Khuri motivates the full-cut quantity
used here \cite{BrayKhuri2010}. This formulation differs from bounds using
the apparent-horizon area itself or the area of an outermost generalized
apparent horizon, for which counterexamples are known
\cite{BenDov2004,CarrascoMars2010}.

Charge changes both the inequality and its equality model. If
$Q^2=Q_E^2+Q_B^2$, the outer Reissner--Nordstr\"om radius is
$m+\sqrt{m^2-Q^2}$. The appropriate statement for possibly disconnected
horizons is the upper-area inequality
\begin{equation}\label{eq:intro-upper-area}
 m\ge Q,\qquad
 \sqrt{\frac{A}{4\pi}}\le m+\sqrt{m^2-Q^2}.
\end{equation}
Writing $r_A=\sqrt{A/(4\pi)}$, the familiar polynomial expression
$m\ge(r_A+Q^2/r_A)/2$ is the required branch only for $r_A>Q$.
Khuri, Weinstein and Yamada developed a charged conformal-flow approach
for multiple boundary components \cite{KhuriWeinsteinYamada2017}.
The numerical argument here instead converts electric flux into a
controlled energy decrease and applies the neutral enclosing-area
theorem \cite[Theorem~1.3]{NeutralEnclosingArea2026}.

This paper proves the minimum-enclosing-area formulation of the charged
spacetime Penrose conjecture for the exterior data of
Definition~\ref{def:data}. The result uses the invariant ADM mass and the
area in the original metric. It permits nonzero ADM momentum,
disconnected weakly future outer-trapped boundary, both source-free
electromagnetic fields, and uncharged non-electromagnetic matter obeying
its dominant energy condition (DEC). Only two metric derivatives and one
second-form derivative are controlled at infinity. Neither a positive
area infimum nor a timelike ADM vector is assumed. Theorem~\ref{thm:numerical}
establishes both and proves \eqref{eq:intro-upper-area}.

The equality theorem concerns the entire original initial data.
For a connected, outer area-minimizing, outermost marginal boundary in the
nondegenerate range $m>Q$, equality holds precisely in the admissible
Reissner--Nordstr\"om class specified in
Theorems~\ref{thm:rigidity} and \ref{thm:converse}. The original metric,
second form and both electromagnetic fields are recovered by a global
spacelike embedding that is smooth through the future horizon. This
includes non-time-symmetric and boosted equality examples. No statement
is made about the region behind the initial boundary.

\subsection{Construction of the numerical comparison}

The proof first prepares a strict exterior with compactly supported second
form and a smooth rest end. The preparation preserves the two closed flux
forms and compares the original enclosing infima and invariant masses.
Its annular construction uses the stated weak derivatives and integrable
constraint sources; it does not impose parity or higher-order decay on
the original data. An energy-raising conformal family then deals with the
possible non-timelike ADM endpoint after the timelike estimate is proved.

On each prepared exterior we solve a coupled scalar system on an
artificial compact filling. One variable is a graph height; the other
controls conformal distortion and compression. A scalar identity provides
positive quadratic terms strong enough to preserve the charged
curvature lower bound after projection of the two flux forms. The system
is solved with its floor, boundary-slope bounds and end normalization,
rather than subject to a separate solvability hypothesis.

Two features of this construction are useful beyond the final
inequality. First, the regularity argument exploits the rank-one axis
structure of a uniformly elliptic equation whose axial coefficient need
not have a modulus of continuity. A trace-reversed Hessian flux and a
fixed-axis comparison give gradient control without differentiating that
coefficient. Second, the estimates distinguish constants at a fixed
approximation parameter from the quantitative distortion bounds needed
when that parameter tends to infinity. The height separates the filling
from a region with a charged scalar-curvature bound. A second bounded
scalar solve converts the total charge into an exact decrease of ADM
energy, while its conformal floor controls every full enclosing area.
A compactly supported pure-trace second form makes a regular height
boundary future trapped, so the neutral theorem applies. Optimizing the
flux profile gives both branches of the charged bound. The Maxwell
forms need be closed only on the original exterior; their smooth
extensions into the artificial filling impose no flux constraint there.

\subsection{Recovery of equality data}

After the numerical inequality, the proof returns to the original
exterior. Variation of the enclosing infimum is performed over its whole
compact family of perimeter minimizers. Constraint separation produces a
causal lapse--shift pair. Variations of both closed flux forms yield exact
electromagnetic Killing contractions even when the exterior is not
simply connected. The resulting stationary data are the original data;
their possible extra stress is an aligned null tensor.

The same variation argument excludes larger marginal obstacles with a
fixed future or past sign on each component. This stronger exclusion
supplies causal communication with the end. A smooth bifurcation
attachment, including an all-orders normal-jet argument at vanishing
lapse, then places the stationary exterior within the hypotheses of the
maximal-hypersurface theorem of Chru\'sciel and Wald
\cite{ChruscielWald1994}. On that maximal slice, the magnetic potential and
momentum identities adapt the integral staticity method of Sudarsky and
Wald \cite{SudarskyWald1993}, with the residual-matter term treated locally,
to force staticity and the disappearance of additional matter. The static
orbit metric has the original invariant mass and boundary area.
The complete electrostatic equations and normalized lapse permit
application of the connected-horizon uniqueness theorem of Borghini,
Cederbaum, and Cogo \cite[Theorem~3.1]{BorghiniCederbaumCogo2025}.
It identifies the static metric, lapse, and electric potential with
Reissner--Nordstr\"om, after which the original slice, horizon embedding
and two signed charges are recovered.

\subsection{Organization}

Section~\ref{sec:statements} gives the full hypotheses and conventions.
The geometric preliminaries justify positivity and the full-frontier
perimeter formulation. The end preparation, filled-system identity,
a priori estimates and elliptic solve construct the comparison data;
the charged comparison closes the numerical proof. The final three
sections carry out original-data variation, the stationary extension and
the rigidity and sharpness arguments.
